\section{Introducción}

El contactor o conmutador electromagnético es un dispositivo de uso industrial cuya función es establecer o interrumpir el paso de corriente en un circuito de potencia, mediante un circuito de magnetización que acciona sus contactos. Se compone de dos partes esenciales: el circuito de magnetización (bobina de control) y el circuito actuador o de potencia (contactos principales y auxiliares). En esta práctica se estudiará el circuito magnético de un contactor comercial, analizando su comportamiento en los estados de reposo y energizado~\cite{mit,schneider}.

\section{Objetivos}

\subsection*{Objetivo General}
\begin{itemize}
    \item Analizar el circuito magnético de un contactor electromagnético comercial, identificando sus componentes y evaluando sus características eléctricas y magnéticas.
\end{itemize}

\subsection*{Objetivos Específicos}
\begin{itemize}
    \item Identificar las partes constitutivas del contactor: bobina, núcleo, armadura o martillo, contactos principales y auxiliares.
    \item Comprender la función de las espiras de sombra (anillos conductores) en el núcleo magnético.
    \item Explicar el propósito del entrehierro permanente entre el núcleo y la armadura.
    \item Definir y calcular la corriente de llamada y la corriente de mantenimiento.
    \item Determinar la tensión mínima de operación y la tensión de mantenimiento.
    \item Comparar los resultados experimentales con los datos de placa del fabricante.
\end{itemize}
